% \xwhatthisisheading\label{sec:what-this-is}

\paragraph{This is a non-technical explainer for active inference as practised} Active inference is normally introduced as mathematics, and the mathematics is
not hard to find.  What is hard to find is an account of what the framework
actually asks you to build --- which decisions it makes for you, which it leaves
open, where it stops being a theory and starts being a set of engineering
choices someone had to make on a Tuesday.  This paper is an attempt at that
account, written for readers who want to understand the shape of the thing
without first working through the formalism.  There is no mathematics in the
sections that follow.  Where a quantity matters, it is named and described in
words.

What distinguishes this from a tutorial is that it describes a system that
exists.  The implemented catalogue entries describe running code, not
sketches presented as completed designs.  The entries were written and revised
\emph{as} the system was built and broken and repaired, and several carry the
scar of the specific failure that forced them.  Entries chartered but not yet
built are explicitly marked as staged rather than allowed to borrow the standing
of implemented entries.  Where a measurement came out against us, the
number is reported anyway.  We take this to be the paper's main contribution to
the pattern community: not the individual moves, most of which will look
familiar to anyone who has built a control system, but the fact that the
catalogue is answerable to a reference implementation and can therefore be
checked, contradicted, and reused rather than merely admired.\worknote{
CHECKED, literally: every operational sentence was audited into a 150-row
inventory --- VALIDATED 14, EXISTS-UNVALIDATED 48, ABSENT-OR-LOST 21,
REFUTED 7, UNCLASSIFIED 60. The seven REFUTED rows await their evidenced
rewrite (row 29); these notes flag affected paragraphs.}

``Patterns'' in the title names the pattern-language subject under
examination, not a status verdict over the catalogue: an entry earns pattern
status through repeated witnessed enactment and review, and no entrywise
census of that kind has been completed.\worknote{UNCHANGED and still true:
no entrywise promotion census exists. The staged mechanism for one ---
pattern-by-signature ticket edges in the apparatus ledger --- is built but
not yet run to completion.}

\paragraph{Who benefits, and how.}  Three audiences, wanting different things:

\begin{itemize}\tightlist
\item \emph{People building agent harnesses.}  The catalogue is usable as a
  checklist for making an agent system answerable: what to record, what to keep
  separate from what, where to refuse to let the system certify itself, and ---
  stated explicitly for each pattern --- the one invariant a different
  implementation would have to preserve.  Most of these patterns are not
  specific to active inference, and several were learned the expensive way.
\item \emph{The pattern community.}  This is a worked case of design patterns
  used \emph{operationally} rather than descriptively: the same pattern text
  that a person reads is also the unit an agent selects, composes, and is held
  to.  That dual use --- readable warrant and machine-selectable action --- is a
  live question for what pattern languages are for in an era of agentic tools,
  and this paper is a data point rather than an argument.
\item \emph{Researchers in active inference and agentic software engineering.}
  The formal treatment --- the generative model, the scoring quantities, the
  structure-learning operator, the faithfulness audit of where our
  implementation departs from the theory, and the preregistered experiment with
  its negative result --- is in the companion paper, which shares this paper's
  system and specimen run.
\end{itemize}
\worknote{RULED (2026-09-13): the companion is futon-2026, which owns the
formal and preregistered experimental material --- the blanket split:
this paper's blanket is the REPL, futon-2026's closure includes the
operator. RESOLVED (2026-09-14): the ruled Methods section
(\S\ref{sec:methods}) now carries the Lean certificate --- the F1--F4
closures over the pinned find record plus its verified transcription
chain. Inventory rows P001--P005 close.}

\paragraph{What this paper does not claim.}  It does not claim that governing
coding agents this way makes them better.  That comparison has not been run: it
would need an arm in which the same agents work the same codebase without
active-inference-governed selection, adjudicated by someone with no stake in the
outcome, over enough attempts to mean something.  What can be shown, and is
shown here, is narrower and still worth having: that the loop closes, that its
decisions carry reasons a person can inspect and disagree with, and that its
claims about the world are backed by evidence it did not manufacture.  A reader
who finishes this paper believing we have demonstrated benefit has been misled,
and we would rather say so at the start than be read that way.\worknote{
UNCHANGED by design: the benefit comparison has still not been run, and
nothing in the 2026-09 campaign changes that. What the campaign did
strengthen is the narrower claim set: ``the loop closes'' is now
demonstrated on demand (run r1), and ``claims backed by evidence it did not
manufacture'' is partially machine-checked by the witness registry.}
